\begin{figure}[!ht]
    \centering
    \includegraphics[width=\textwidth]{figs/overview2.pdf}
    \vspace{-5 mm}
    \caption{\textbf{Overview of DynFlowDrive.} Given current observations, multi-mode trajectories are firstly generated by the standard planning module. A flow-based dynamic latent world model is incorporated to simulate the progressive future evolution in latent space. The resulting dynamics are used by a stability-aware multi-mode selection module, which assess the trajectory based on reconstruction quality and flow-based stability, enabling reliable supervision and improved planning robustness.}
    \vspace{-3mm}
    \label{fig:overview} %
\end{figure}

\section{Method}


As illustrated in Figure~\ref{fig:overview}, DynFlowDrive is built upon a standard end-to-end planning framework together with a flow-based latent world modeling. Given the surrounding $N$-view camera images 
$\mathbf{I}_t = \{ I_t^i \}_{i=1}^N$ and a high-level navigation command $cmd$ at current timestep $t$, a vision-based approach aims to predict the planning trajectory 
$\mathbf{T}$ that consists of a set of future points in BEV space. 

\noindent DynFlowDrive comprises three key parts:
\textbf{1) Trajectory Planning Paradigm} (Sec.~\ref{origin}), which produces multi-mode trajectories based on the observations following the standard planning paradigm in end-to-end autonomous driving;
\textbf{2) Dynamic World Model Prediction} (Sec.~\ref{dy_world}), that simulates latent scenes' evolution conditioned on the planning trajectory with the flow-based dynamic model, and 
\textbf{3) Stability-aware Multi-mode Selection} (Sec.~\ref{mode_select}), which selects trajectory modes with respect to the dynamic stability for reliable supervision. These modules work in a complementary manner, coupling dynamic world simulation with multi-mode trajectory selection.

\subsection{Trajectory Planning Paradigm}
\label{origin}
As shown in the yellow part of Figure~\ref{fig:overview}, the trajectory planning paradigm predicts multiple candidate trajectories. It first encodes the sensor inputs into a unified scene representation, %parameterized by the scene queries. Based on the initialized multi-mode trajectory anchors, 
a trajectory decoder then attends to the scene queries to iteratively refine the trajectory queries into final trajectory proposals.

\noindent{\textbf{Scene Encoder.}} Given multi-view images $\mathbf{I}_t=\{I_t^i\}_{i=1}^V$ at the current timestep $t$, each view is encoded into perspective-view features $\mathbf{F}_t^{\text{pv}} = \{\mathbf{F}_t^i\}_{i=1}^V$. A set of learnable scene queries $\mathbf{Q}^s_t$ is then introduced to aggregate multi-view information via cross-attention, which is formulated as:
\begin{equation}
\mathbf{Q}_t^{\text{scene}} = \mathrm{CrossAttn}(\mathbf{Q}_t^{\text{scene}}, \mathbf{F}_t^{\text{i}}).
\end{equation}
It yields a unified scene representation that captures the spatial layout and context from the driving scenes of the current time step.

\noindent{\textbf{Multi-modal Trajectory Decoder.}} Since $\mathbf{Q}_t^{\text{scene}}$ contains relevant scene information, following~\cite{jiang2023vad, hu2023planning}, we initialize a set of candidate waypoint queries $\mathbf{Q}_t^{\text{traj}}=\{q_t^i\}_{n=1}^{N} \in \mathbb{R}^{N_m\times N \times \times D}$ to extract multi-mode planning trajectories, where $N_m$ denotes the number of driving commands, $N$ denotes the number of trajectory modes, and $D$ is the dimensions of query feature. %Specifically, we introduce a set of $N$ learnable trajectory queries , where each query corresponds to one trajectory mode and serves as an anchor to be refined.

Then, scene-level context $\mathbf{Q}_t^{\text{scene}}$ is incorporated into the trajectory queries through cross-attention with the scene queries:
\begin{equation}
\mathbf{Q}_t^{\text{traj}} = \mathrm{CrossAttn}(\mathbf{Q}_t^{\text{traj}}, \mathbf{Q}_t^{\text{scene}}).
\end{equation}

Following the context-enhanced trajetcory queries, two diiferent MLP heads are applied for trajectory refinement and mode scoring. The trajectory head $\mathrm{MLP}_{\text{traj}}$ predicts a residual update to refine the waypoints, while the score head $\mathrm{MLP}_{\text{s}}$ estimates a confidence score for each mode, which can be formulated as:
\begin{equation}
\hat{\mathbf{T}}_t = \mathbf{T}_t + \mathrm{MLP}_{\text{traj}}(\mathbf{Q}_t^{\text{traj}}), \quad
\mathbf{c}_t = \mathrm{MLP}_{\text{s}}(\mathbf{Q}_t^{\text{traj}}),
\end{equation}
where $\hat{\mathbf{T}}_t=\{\hat{T}_t^n\}_{n=1}^{N}$ denotes the refined trajectories, and $\mathbf{c}_t=\{c_t^n\}_{n=1}^{N}$ represents their corresponding mode scores. These $(\hat{\mathbf{T}}_t,\mathbf{c}_t)$ pairs are treated as the multi-modal trajectory proposals for subsequent selection and training supervision.

%Given the ground-truth trajectory $\mathbf{T}^{\mathrm{gt}}=\{\mathbf{p}^{\mathrm{gt}}_l\}_{l=1}^{L}$, we supervise the predicted trajectories using an $\ell_2$ regression loss:
%\begin{equation}
%\mathcal{L}_{\mathrm{traj}}=\frac{1}{L}\sum_{l=1}^{L}\left\|\hat{\mathbf{p}}_{l}-\mathbf{p}^{\mathrm{gt}}_{l}\right\|_2^2,
%\end{equation}
%where 

This formulation provides a set of diverse trajectory hypotheses while maintaining a unified query-based decoding structure, which serves as the foundation for subsequent trajectory selection and dynamic world modeling.

\subsection{Dynamic Latent World Model}
\label{dy_world}

To explicitly characterize planning quality through trajectory-induced world evolution, we introduce a dynamic latent world model that simulates how the environment evolves under different trajectory hypotheses. Given the current latent state and candidate trajectories, the model predicts future latent states via, which are used to evaluate the resulting trajectory outputs in latent space.

\noindent{\textbf{World Feature Extraction.}}
For the dynamic world simulator, directly reusing the driving encoder may lead to representation shifts and unstable features for modeling temporal evolution.  Rather than follow prior works~\cite{li2024law,li2024ssr}, we employ a pretrained foundation encoder to obtain more stable and dynamics-aligned latent representations.

Given the surrounding $V$-view images $\mathbf{I}_t = \{ I_t^i \}_{i=1}^V$, we firstly encode each view into a latent representation using a pretrained VAE encoder $\mathbf{E}_{\text{vae}}$ followed by the lightweight MLP :
\begin{equation}
\mathbf{z}_t^i = \mathrm{MLP}(\mathbf{E}_{\text{vae}}(I_t^i)), \quad i = 1, \dots, V.
\end{equation}
where $\tilde{\mathbf{z}}_t^i$ denotes the dynamics-aware latent feature for the $i$-th view. Such representation benefits from the diverse training data and strong generalization capability of foundation models to obtain dynamics-oriented features.



To incorporate global scene contexts, we aggregate the multi-view latent features by interacting them with the scene queries $\mathbf{Q}_t^{\text{scene}}$ via cross-attention:
\small{
\begin{equation}
\tilde{\mathbf{z}}_t^{\text{w}} = \mathrm{CrossAttn}(\mathbf{Q}_t^{\text{scene}}, \{\mathbf{z}_t^i\}_{i=1}^V),
\end{equation}}
where $\tilde{\mathbf{z}}_t^{\text{w}}$ denotes the world latent that captures both multi-view spatial features and global scene context.

Following the same procedure, we extract the corresponding world features for the next timestep, denoted as $\tilde{\mathbf{z}}^{\text{w}}_{t+1}$, which serve as the supervision target for training the dynamic flow model.


\begin{wrapfigure}{r}{0.66\textwidth}
    \centering
    \vspace{-24pt}
\includegraphics[width=0.66\textwidth]{figs/wmdesign.pdf}
    \vspace{-16pt}
    \caption{The architecture of our \textbf{dynamic latent world model design}, in which the velocity field $v_{\theta}$ is learnt to capture the trajectory-conditioned dynamics transitions in the latent space. }
    \vspace{-20pt}
    \label{fig:wm-design}
\end{wrapfigure}

\noindent{\textbf{Flow-based Latent World Model Simulation.}}
Given the current world latent $\tilde{\mathbf{z}}^{\text{w}}_t$ and predicted trajectories $\hat{\mathbf{T}}_t$, we model trajectory-conditioned world evolution in latent space. As shown in Figure~\ref{fig:wm-design}, instead of learning a generic mapping between noise and data, we aim to learn a trajectory-conditioned velocity field that describes how the latent state evolves under different motion hypotheses.
Following the rectified flow formulation~\cite{liu2022flow},
 %we learn a velocity field $v_\theta(\cdot)$ that transports the current latent toward its future state by a transformer-based structure. Given the target latent $\mathbf{z}^{\text{w}}_{t+1}$,, and add stochastic perturbations:
 we construct the noised anchor state $\mathbf{a}$ from the current latent to provide a stochastic starting point for modeling local state transitions:
 \begin{equation}
\mathbf{a} = (1-\alpha)\tilde{\mathbf{z}}_t^{\text{w}} + \alpha \boldsymbol{\epsilon}, \quad 
\boldsymbol{\epsilon} \sim \mathcal{N}(\mathbf{0}, \mathbf{I}), 
\end{equation}
where $\boldsymbol{\epsilon}$ is the Gaussian noise with the same shape as $\mathbf{z}_t^{\text{w}}$, and $\alpha \in [0,1]$ controls the perturbation strength.
Then the continuous interpolation between the anchor $\mathbf{a}$ and the target latent $\tilde{\mathbf{z}}_{t+1}^{\text{w}}$ is then defined as:
\begin{equation}
\mathbf{x}_s = (1-s)\mathbf{a} + s\,\tilde{\mathbf{z}}_{t+1}^{\text{w}}, \quad s \sim \mathcal{U}(0,1),
\end{equation}
where $s$ denotes the flow timestep.


To incorporate the trajectory conditions, planning modes $\{\hat{T}_t^n\}_{i=n}^N$ are encoded into a trajectory embeddings by $\mathrm{TrajEmb}(\cdot)$ and fused with the original VAE latent to obtain the condition features $\{\mathbf{h}_t^n\}_{n=1}^{N}$:
\begin{equation}
\mathbf{h}_t^n = \mathrm{Concat}\big([\lambda_z\cdot \mathbf{z}^{\text{w}}_t,\ \lambda_T \cdot\mathrm{TrajEmb}(\hat{T}_t^n)]\big),
\end{equation}
where $\lambda_z$ and $\lambda_T$ are scalar weights to balance the latent state and trajectory embeddings.
Then, the velocity field can be predicted by a transformer-based flow model $\mathcal{F}_\theta(\cdot)$, which is formulated as:
\begin{equation}
v_\theta = \mathcal{F}_\theta(\mathbf{x}_s, s, \mathbf{h}_t^i).
\end{equation}
This formulation models how the latent state evolves under the corresponding trajectory hypothesis.

Thus, during training, we adopt the flow matching objective that supervises the velocity along the interpolation path between the anchor state $\mathbf{a}$ and the target latent $\tilde{\mathbf{z}}_{t+1}^{\text{w}}$:
\begin{equation}
\mathcal{L}_{flow} =
\mathbb{E}
\left[
\left\|
\mathcal{F}_\theta(\mathbf{x}_s, s, \mathbf{h}_t^{i})
-
(1-s)(\tilde{\mathbf{z}}_{t+1}^{\text{w}}-\mathbf{a})
\right\|_2^2
\right],
\label{eq_loss}
\end{equation}
where the target velocity corresponds to the displacement from the anchor state toward the future latent along the interpolation path.

\noindent For sampling integration, the learned velocity defines a trajectory-conditioned dynamical system so that the future world latent can be obtained via integration:
\begin{equation}
\frac{d\tilde{\mathbf{z}}^{\text{w}}}{ds}
=
\mathcal{F}_\theta(\tilde{\mathbf{z}}^{\text{w}}(s), s, \mathbf{h}_t^{i}).
\end{equation}
Each step corresponds to an integration step along the flow variable $s$, which simulates a small portion of the latent evolution between the driving timesteps $t$ and $t+1$.

Thus, starting from the current latent $\tilde{\mathbf{z}}_t^{\text{w}}$, the future world latent can be obtained by integrating the velocity field from $s=0$ to $1$. 
In practice, we approximate this integration using a discrete Euler solver:
\begin{equation}
\tilde{\mathbf{z}}_{s_{k+1}}^{\text{w}}
=
\tilde{\mathbf{z}}_{s_k}^{\text{w}}
+
\Delta s \,
\mathcal{F}_\theta(\tilde{\mathbf{z}}_{s_k}^{\text{w}}, s_k, \mathbf{h}_t^{i}),
\end{equation}
where $s_k$ is the corresponding flow timestep, and $\Delta s$ is the integration step size.


This formulation models how different trajectory hypotheses induce distinct latent state transitions, enabling dynamics-aware trajectory evaluation based on the progressive evolution dynamics.


\begin{wrapfigure}{r}{0.35\textwidth}
    \centering
    \vspace{-56pt}
\includegraphics[width=0.35\textwidth]{figs/selector.pdf}
    \vspace{-16pt}
    \caption{Stability-aware Multi-mode Selection. For training, the score head is supervised by the stable criterion. For Inference, the best mode trajectory is selected according to the highest score index.} 
    \vspace{-26pt}
    \label{fig:selector}
\end{wrapfigure}

%\subsection{Stable-aware Multi-mode Selection}
\subsection{Stability-aware Multi-mode Selection}
\label{mode_select}



Our multi-modal trajectory predictor generates a set of candidate trajectories 
$\{\hat{\mathbf{T}}_t^{n}\}_{n=1}^{N}$ to capture the inherent uncertainty of future driving behaviors. 
Conventional selection strategies typically rely on geometric or reconstruction criteria, 
such as trajectory error (e.g., ADE/FDE) and latent reconstruction loss. 
However, these metrics only measure spatial or feature-level similarity and fail to reflect the quality of trajectory-induced world evolution.

In practice, different trajectory candidates may produce similar geometric errors while leading to substantially different latent state evolutions. 
Selecting trajectories solely based on geometric criteria may therefore favor modes that are spatially accurate but induce unstable or unrealistic scene dynamics, which is undesirable for reliable planning.

To address this issue, as shown in Figure~\ref{fig:selector}, we introduce a stability-aware selection strategy that evaluates each trajectory based on the induced dynamics.
Given a trajectory candidate $\hat{\mathbf{T}}_t^{n}$, the flow-based world model simulates the latent evolution conditioned on this trajectory and produces a sequence of velocities 
$\{\mathbf{v}_k^{n}\}_{k=1}^{K}$ along the integration steps, where $K$ is the discrete steps during transitions. 
We measure the directional consistency of the evolution by computing the average angular deviation between consecutive normalized velocity directions $\hat{\mathbf{u}}_{k-1}^{n}$ and $\hat{\mathbf{u}}_{k}^{n}$:
\begin{equation}
\hat{\mathbf{u}}_k^{n} = \frac{\mathbf{v}_k^{n}}{\|\mathbf{v}_k^{n}\|}, \quad
\mathcal{S}_{\theta}^{n} = \frac{1}{K-1}\sum_{k=2}^{K}\arccos\!\left(\hat{\mathbf{u}}_k^{n} \cdot \hat{\mathbf{u}}_{k-1}^{n}\right),
\end{equation}
where smaller $\mathcal{S}_{\theta}^{n}$ value indicates smoother and more physically consistent dynamics.

Based on this, we define the unified selection criterion, combined with other commonly used metrics~\cite {zheng2025world4drive} for each trajectory mode assessment:
\begin{equation}
\mathcal{C}^{n} = 
\lambda_{rec} \, \mathcal{L}_{rec}^{n} 
+ \lambda_{traj} \, \mathcal{L}_{traj}^{n} 
+ \lambda_{\theta} \, \mathcal{S}_{\theta}^{n}, 
\quad
n^* = \arg\min_n \mathcal{C}^{n},
\end{equation}
where $\mathcal{L}_{traj}^{n}$ denotes the trajectory error between the predicted trajectory and the ground truth, and $\mathcal{L}_{rec}^{n}$ measures the reconstruction discrepancy of the predicted latent state.
The selected optimal mode $n^*$ is employed to supervise the learning of the trajectory scoring head.

By explicitly incorporating trajectory-conditioned dynamical stability, the proposed strategy favors trajectories that induce smooth and physically consistent world evolution, leading to more reliable trajectory assessment beyond conventional criteria.

\subsection{Training and Inference Settings}

For model training, we jointly optimize trajectory prediction, scoring, latent reconstruction, and flow-based dynamics with a unified objective:
\begin{equation}
\mathcal{L} = 
\mathcal{L}_{traj} 
+ \lambda_{score}\mathcal{L}_{score} 
+ \lambda_{rec}\mathcal{L}_{rec} 
+ \lambda_{flow}\mathcal{L}_{flow}.
\end{equation}
Here, $\mathcal{L}_{traj}$ denotes the $L_1$ loss between predicted trajectories and ground truth. 
The scoring loss $\mathcal{L}_{score}$ is supervised using the selected optimal mode $n^*$ from Sec.~\ref{mode_select}. 
$\mathcal{L}_{rec}$ enforces consistency between predicted and target latent states, and $\mathcal{L}_{flow}$ corresponds to the flow matching objective for velocity learning. 
$\lambda_{score}$, $\lambda_{rec}$, and $\lambda_{flow}$ are balancing weights.

Here, $\mathcal{L}_{traj}$ denotes the $L_1$ loss between predicted trajectories and the ground-truth future trajectory. 
The scoring loss $\mathcal{L}_{score}$ is supervised using the selected optimal mode $n^*$ defined in Sec.~\ref{mode_select}. 
$\mathcal{L}_{rec}$ enforces consistency between the predicted latent state and the target latent representation, while $\mathcal{L}_{flow}$ corresponds to the flow matching objective used to learn the trajectory-conditioned velocity field. 
The coefficients $\lambda_{score}$, $\lambda_{rec}$, and $\lambda_{flow}$ balance the contributions of different objectives.

During inference, the world model is not involved, since the driving system has already acquired the capability to select the optimal trajectory. 
The planning module directly selects the trajectory with the highest predicted score as the final output, introducing no additional computational overhead compared with the standard trajectory prediction framework.

